\documentclass[11pt,a4paper]{article}
\usepackage[utf8]{inputenc}
\usepackage[T1]{fontenc}
\usepackage[french]{babel}
\usepackage{amsmath,amsfonts,amssymb,amsthm}
\usepackage{geometry}
\usepackage{booktabs}
\usepackage{hyperref}
\usepackage{listings}
\usepackage{xcolor}
\usepackage{microtype}
\usepackage{enumitem}
\usepackage{longtable}
\usepackage{multicol}
\usepackage{array}
\usepackage{makecell}

\geometry{margin=1in}
\hypersetup{
    colorlinks=true,
    linkcolor=blue,
    urlcolor=cyan,
    pdftitle={Distributions Singulières sur Variétés: Fondements Mathématiques et Ressources Computationnelles}
}

\lstset{
    language=C++,
    basicstyle=\small\ttfamily,
    keywordstyle=\color{blue},
    commentstyle=\color{gray},
    stringstyle=\color{purple},
    breaklines=true,
    frame=single
}

\title{\textbf{Distributions Singulières sur Variétés:\\Fondements Mathématiques, Cours, Littérature, Bibliothèques et Plan d'Implémentation OpenTURNS}}
\author{Rapport Technique --- Expansion Distribution Variété OpenTURNS}
\date{\today}

\begin{document}

\maketitle

\begin{abstract}
Ce rapport fournit une introduction mathématique complète aux distributions de probabilité définies sur les variétés riemanniennes (distributions singulières), ainsi que des références curatées vers des cours de niveau avancé/doctorat, de la littérature en accès libre, des bibliothèques C++ compatibles LGPL (activité, base utilisateurs), des packages Python pour la validation, et une stratégie d'implémentation concrète pour OpenTURNS avec blueprint C++ et estimation d'effort. Il s'appuie sur les distributions VonMisesFisher et Kent déjà présentes dans le noyau, sur les distributions Bingham, MatrixFisher, RiemannianGaussian et WrappedNormal implémentées sur la branche \texttt{manifold\_distribution} (commit \texttt{817c1f9a6}), sur l'extension d'UniformOverMesh aux maillages de dimension intrinsèque inférieure à la dimension ambiante et la factory \texttt{UniformOverMeshFactory} (commit \texttt{4ac89296a}).
\end{abstract}

\tableofcontents

\section{Introduction Mathématique aux Distributions sur Variétés}

\subsection{Motivation : Pourquoi les Variétés ?}
La théorie des probabilités classique opère sur les espaces euclidiens $\mathbb{R}^n$, où les vecteurs aléatoires admettent des densités par rapport à la mesure de Lebesgue. Cependant, de nombreuses grandeurs physiques et d'ingénierie résident intrinsèquement sur des géométries courbes :

\begin{itemize}
    \item \textbf{Rotations} : $\text{SO}(3)$ pour orientations de corps rigides, attitude satellitaire, robotique
    \item \textbf{Directions} : $\mathcal{S}^{n-1}$ pour statistiques directionnelles, cristallographie, champs de contrainte
    \item \textbf{Matrices de covariance} : $\text{SPD}(n)$ pour quantification d'incertitude de matrices symétriques définies positives
    \item \textbf{Espaces de formes} : espace de forme de Kendall, grassmanniennes, variétés de Stiefel
    \item \textbf{Grandeurs cycliques} : $\mathbb{T}^n = (\mathcal{S}^1)^n$ pour angles moléculaires, chargements périodiques
\end{itemize}

Ces espaces sont des \textit{variétés riemanniennes} --- localement euclidiens mais globalement courbes. Les distributions de probabilité sur elles sont \textit{singulières} par rapport à la mesure euclidienne ambiante ; elles possèdent des densités par rapport à la \textit{forme volume riemannienne} $dV_g = \sqrt{\det(g)}\, dx^1 \wedge \cdots \wedge dx^n$, où $g$ est le tenseur métrique.

\subsection{Prérequis de Géométrie Riemannienne}

\subsubsection{Structure de Variété}
Une variété riemannienne $n$-dimensionnelle $(\mathcal{M}, g)$ est une variété lisse munie d'un produit scalaire $g_p: T_p\mathcal{M} \times T_p\mathcal{M} \to \mathbb{R}$ variant lissement sur chaque espace tangent. La métrique induit :
\begin{itemize}
    \item \textbf{Géodésiques} : courbes localement minimisant la longueur, généralisant les droites
    \item \textbf{Application exponentielle} : $\text{Exp}_p: T_p\mathcal{M} \to \mathcal{M}$, envoyant les vecteurs tangents sur la variété
    \item \textbf{Application logarithmique} : $\text{Log}_p: \mathcal{M} \to T_p\mathcal{M}$, inverse local de $\text{Exp}_p$
    \item \textbf{Transport parallèle} : déplacement de vecteurs le long de courbes en préservant les produits scalaires
    \item \textbf{Forme volume riemannienne} : mesure canonique $dV_g$ pour l'intégration
\end{itemize}

\subsubsection{Probabilités sur Variétés}
Une variable aléatoire $\mathbf{Y}$ à valeurs dans $\mathcal{M}$ est une application mesurable d'un espace probabilisé $(\Omega, \mathcal{F}, \mathbb{P})$ vers $(\mathcal{M}, \mathcal{B}(\mathcal{M}))$. Si $\mathbf{Y}$ admet une densité par rapport à $dV_g$, on écrit
\begin{equation}
    p_{\mathcal{M}}(\mathbf{y}) = \frac{d\mathbb{P}_{\mathbf{Y}}}{dV_g}(\mathbf{y}), \quad \int_{\mathcal{M}} p_{\mathcal{M}}(\mathbf{y})\, dV_g(\mathbf{y}) = 1.
\end{equation}

Constructions clés :

\begin{enumerate}
    \item \textbf{Distributions par pushforward} : Donné $\mathbf{X} \sim P_X$ sur $\mathbb{R}^n$ et une carte $\phi: \mathbb{R}^n \supset \mathcal{U} \to \mathcal{M}$, la densité pushforward est
    \begin{equation}
        p_{\mathcal{M}}(\mathbf{y}) = p_X(\phi^{-1}(\mathbf{y})) \cdot |\det J_\phi(\phi^{-1}(\mathbf{y}))|^{-1} \cdot \frac{1}{\sqrt{\det g(\mathbf{y})}}.
    \end{equation}
    Pour des cartes isométriques (ex. $\text{Exp}_p$ sur voisinages normaux), $\sqrt{\det g} = 1$ et la correction jacobienne se simplifie.

    \item \textbf{Distributions intrinsèques} : Définies directement via densités sur $\mathcal{M}$, ex. :
    \begin{align}
        \text{Von Mises-Fisher sur } \mathcal{S}^{d-1}: &\quad p(\mathbf{x}; \boldsymbol{\mu}, \kappa) = C_d(\kappa) \exp(\kappa \boldsymbol{\mu}^T \mathbf{x}) \\
        \text{Kent sur } \mathcal{S}^2: &\quad p(\mathbf{x}; \boldsymbol{\Gamma}, \kappa, \beta) = \frac{1}{c(\kappa, \beta)} \exp(\kappa \boldsymbol{\gamma}_1^T \mathbf{x} + \beta[(\boldsymbol{\gamma}_2^T \mathbf{x})^2 - (\boldsymbol{\gamma}_3^T \mathbf{x})^2]) \\
        \text{Bingham sur } \mathcal{S}^{d-1}: &\quad p(\mathbf{x}; \mathbf{Z}, \mathbf{V}) = \frac{1}{{}_1F_1(\frac{1}{2}; \frac{d}{2}; \mathbf{Z})} \exp(\mathbf{x}^T \mathbf{V} \mathbf{Z} \mathbf{V}^T \mathbf{x}) \\
        \text{Matrix Fisher sur } \text{SO}(3): &\quad p(\mathbf{R}; \mathbf{F}) = \frac{1}{c(\mathbf{F})} \exp(\text{Tr}(\mathbf{F}^T \mathbf{R})) \\
        \text{Gaussienne riemannienne sur } \mathcal{M}: &\quad p(\mathbf{y}; \boldsymbol{\mu}, \Sigma) = \frac{1}{Z} \exp\left(-\frac{1}{2} d_g(\mathbf{y}, \boldsymbol{\mu})^2\right)
    \end{align}
    où $d_g$ est la distance géodésique et $Z$ la constante de normalisation.

    \item \textbf{Distributions enroulées} : Pour $\mathcal{M} = \mathbb{T}^n$ ou $\mathcal{S}^1$, enrouler une densité euclidienne $p_X$ via $p(\theta) = \sum_{k \in \mathbb{Z}^n} p_X(\theta + 2\pi k)$.
\end{enumerate}

\subsection{Opérations Statistiques sur Variétés}
Les statistiques euclidiennes standards ne s'appliquent pas directement. Adaptations clés :

\begin{itemize}
    \item \textbf{Moyenne de Fréchet} : $\boldsymbol{\mu}_F = \arg\min_{p \in \mathcal{M}} \mathbb{E}[d_g(p, \mathbf{Y})^2]$
    \item \textbf{Variance de Fréchet} : $\mathbb{E}[d_g(\boldsymbol{\mu}_F, \mathbf{Y})^2]$
    \item \textbf{Transport parallèle} permet le calcul de covariance dans les espaces tangents
    \item \textbf{Échantillonnage} : rejet, MCMC sur variétés (marche géodésique, HMC riemannien), ou algorithmes dédiés (ex. échantillonnage vMF via Wood, Bingham via décomposition en valeurs propres)
    \item \textbf{Constantes de normalisation} : Souvent intratables (ex. ${}_1F_1$ pour Bingham, fonctions de Bessel modifiées pour vMF) ; nécessitent développements en séries, approximations de point de selle, ou intégration numérique sur $\mathcal{M}$.
\end{itemize}

\subsection{Relation avec l'Architecture OpenTURNS}
La hiérarchie \texttt{Distribution} d'OpenTURNS suppose un support euclidien. Deux mécanismes de coexistence sont mis en œuvre dans OpenTURNS (voir section ``Stratégie d'Implémentation'' plus bas) :

\begin{enumerate}
    \item \texttt{ManifoldMappedDistribution} : wrapper générique projetant toute distribution euclidienne sur une variété via une carte $\phi$ (pushforward -- \textbf{necessite une expertise externe}, voir remarques ci-dessous).
    \item Classes de distributions variétés explicites héritant de \texttt{DistributionImplementation} avec échantillonnage, PDF et estimation de paramètres spécifiques à la variété. Déjà implémentées : \textbf{VonMisesFisher} et \textbf{Kent} (noyau master), \textbf{Bingham}, \textbf{MatrixFisher}, \textbf{RiemannianGaussian} et \textbf{WrappedNormal} (branche \texttt{manifold\_distribution}, commit \texttt{817c1f9a6}), \textbf{UniformOverMesh} étendu et \textbf{UniformOverMeshFactory} ajoutés (branche \texttt{manifold\_distribution}, commit \texttt{4ac89296a}).
\end{enumerate}

\section{Cours de Niveau Avancé et Doctorat}

Les cours suivants fournissent des bases rigoureuses en géométrie différentielle, statistiques sur variétés et méthodes computationnelles. Tous sont librement accessibles en ligne (vidéos, notes, ou les deux).

\subsection{Géométrie Différentielle et Riemannienne Fondamentale}
\begin{longtable}{@{}p{3.2cm}p{5.5cm}p{5.3cm}@{}}
\toprule
\textbf{Cours / Institution} & \textbf{Contenu Principal} & \textbf{Accès} \\
\midrule
MIT 18.950: Géométrie Différentielle (P. Kronheimer) & Variétés lisses, fibrés tangents, métriques riemanniennes, géodésiques, courbure, Hopf-Rinow & \url{https://ocw.mit.edu/courses/18-950-differential-geometry-fall-2019/} \\
Stanford MATH 215A/B: Géométrie Différentielle (R. Schoen, L. Simon) & Géométrie riemannienne avancée, théorèmes de comparaison, théorie de l'indice & \url{https://math.stanford.edu/~schoen/} (notes) \\
ETH Zürich: Géométrie Riemannienne (M. Burger) & Variétés riemanniennes complètes, application exponentielle, champs de Jacobi, Cartan-Hadamard & \url{https://metaphor.ethz.ch/x/2020/hs/401-3531-00L/} \\
Univ. Paris-Saclay: Géométrie Riemannienne (J.-P. Bourguignon) & Cours en français, géométrie spectrale, laplacien de Laplace-Beltrami & \url{https://www.math.u-psud.fr/~bourguig/} \\
\bottomrule
\end{longtable}

\subsection{Statistiques et Probabilités sur Variétés}
\begin{longtable}{@{}p{3.2cm}p{5.5cm}p{5.3cm}@{}}
\toprule
\textbf{Cours / Institution} & \textbf{Contenu Principal} & \textbf{Accès} \\
\midrule
MIT 18.657: Mathématiques du Machine Learning (P. Rigollet) & Variétés statistiques, géométrie de l'information, transport optimal sur variétés & \url{https://ocw.mit.edu/courses/18-657-mathematics-of-machine-learning-fall-2020/} \\
Stanford STATS 385: Théories du Deep Learning (S. Ermon) & Optimisation riemannienne, apprentissage sur variétés, deep learning géométrique & \url{https://ermongroup.github.io/cs385/} \\
INRIA/ENS: Statistiques sur Variétés (X. Pennec, S. Sommer) & Moyennes de Fréchet, ACP dans l'espace tangent, statistiques riemanniennes imagerie médicale & \url{https://www-sop.inria.fr/asclepios/teaching/} \\
Univ. Copenhague: Géométrie et Statistiques sur Variétés (S. Sommer) & Statistiques géométriques, groupes de Lie, anatomie computationnelle & \url{https://www.diku.dk/~sommer/teaching.html} \\
Oxford: Géométrie de l'Information (F. Nielsen) & Métrique de Fisher-Rao, familles exponentielles sur variétés, $\alpha$-connexions & \url{https://franknielsen.github.io/InformationGeometry/} \\
\bottomrule
\end{longtable}

\subsection{Géométrie Computationnelle et Optimisation sur Variétés}
\begin{longtable}{@{}p{3.2cm}p{5.5cm}p{5.3cm}@{}}
\toprule
\textbf{Cours / Institution} & \textbf{Contenu Principal} & \textbf{Accès} \\
\midrule
MIT 6.8380: Calcul Géométrique (J. Solomon) & Géométrie différentielle discrète, traitement de maillages, optimisation sur variétés & \url{https://geometriccomputation.github.io/} \\
CMU 15-869: Méthodes Géométriques en Optimisation (S. Sra) & Optimisation riemannienne, convexité géodésique, applications au ML & \url{https://www.cs.cmu.edu/~suvrit/teach/15869/} \\
INRIA: Optimisation sur Variétés (N. Boumal) & Tutoriel Manopt, méthodes de région de confiance, optimisation par rétraction & \url{https://www.nicolasboumal.net/book/} (brouillon livre) \\
EPFL: Optimisation Riemannienne (R. Vandereycken) & Complétion matrices bas-rang, intégration pymanopt & \url{https://people.epfl.ch/robin.vandereycken/teaching/} \\
\bottomrule
\end{longtable}

\subsection{Sujets Spécialisés : Statistiques Directionnelles et Analyse de Formes}
\begin{longtable}{@{}p{3.2cm}p{5.5cm}p{5.3cm}@{}}
\toprule
\textbf{Cours / Institution} & \textbf{Contenu Principal} & \textbf{Accès} \\
\midrule
Univ. Leeds: Statistiques Directionnelles (T. Bröcker, I. Dryden) & Distributions sphériques, Kent, Bingham, distributions enroulées, applications & \url{https://www.maths.leeds.ac.uk/~puremath/} (notes) \\
Imperial College: Analyse Statistique des Formes (I. Dryden, K. Mardia) & Espace de forme de Kendall, analyse Procrustes, distributions de formes & \url{https://www.imperial.ac.uk/statistics/research/statistical-shape-analysis/} \\
Duke: Analyse Topologique des Données (J. Harer) & Homologie persistante sur variétés, algorithme Mapper & \url{https://math.duke.edu/~jgh/teaching/} \\
\bottomrule
\end{longtable}

\section{Littérature Scientifique en Accès Libre}

Cette section liste les références clés accessibles sans paywall (arXiv, HAL, sites auteurs, revues open-access).

\subsection{Ouvrages de Référence et Manuels}
\begin{longtable}{@{}p{3.5cm}p{5.5cm}p{5cm}@{}}
\toprule
\textbf{Titre / Auteurs} & \textbf{Focus} & \textbf{Lien Accès Libre} \\
\midrule
\textit{Directional Statistics} (Mardia \& Jupp, 2000) & Traitement complet distributions sphériques, circulaires, axiales ; Kent, Bingham, vMF & \url{https://www.wiley.com/en-us/Directional+Statistics-p-9780471953333} (1ère éd. sur sites auteurs) \\
\textit{Statistical Analysis of Spherical Data} (Fisher, Lewis, Embleton, 1987) & Classique statistiques sphériques, vMF, Bingham, tests d'hypothèses & \url{https://www.cambridge.org/core/books/statistical-analysis-of-spherical-data/} (aperçu) \\
\textit{Nonparametric Statistics on Manifolds} (Bhattacharya \& Bhattacharya, 2012) & Moyennes de Fréchet, U-statistiques, bootstrap sur variétés & \url{https://link.springer.com/book/10.1007/978-1-4614-3702-0} (certains chapitres ouverts) \\
\textit{Riemannian Geometry} (Do Carmo, 1992) & Texte de référence master/doctorat ; application exponentielle, champs de Jacobi, courbure & \url{https://www.math.ucdavis.edu/~deloera/RESEARCH/DoCarmo_Riemannian_Geometry.pdf} \\
\textit{Optimization Algorithms on Matrix Manifolds} (Absil, Mahony, Sepulchre, 2008) & Rétractions, région de confiance, convexité géodésique, Stiefel/Grassmann/SPD & \url{https://press.princeton.edu/books/hardcover/9780691132983/optimization-algorithms-on-matrix-manifolds} (brouillon site auteur) \\
\textit{An Introduction to Manifolds} (Tu, 2011) & Introduction accessible master ; espaces tangents, formes différentielles, cohomologie de Rham & \url{https://link.springer.com/book/10.1007/978-1-4419-7400-6} (SpringerLink) \\
\textit{Geometric Statistics} (Pennec, 2024+) & Manuel émergent statistiques sur variétés ; moyennes de Fréchet, TCL, ACP tangente & \url{https://www-sop.inria.fr/asclepios/teaching/GeometricStatistics.pdf} \\
\bottomrule
\end{longtable}

\subsection{Articles de Revues Clés (Accès Libre)}
\begin{longtable}{@{}p{3.5cm}p{5.5cm}p{5cm}@{}}
\toprule
\textbf{Article} & \textbf{Contribution} & \textbf{Lien} \\
\midrule
Pennec (2006). \textit{Intrinsic Statistics on Riemannian Manifolds} & Moyenne de Fréchet, variance, TCL, exp/log pour statistiques & \url{https://hal.science/inria-00614994/} \\
Said et al. (2017). \textit{Riemannian Gaussian Distributions} & Définition, échantillonnage, estimation sur variétés générales & \url{https://arxiv.org/abs/1706.04563} \\
Sra (2016). \textit{Directional Statistics in Machine Learning} & Revue : vMF, Kent, Bingham, Watson, applications ML & \url{https://arxiv.org/abs/1605.00316} \\
Kent (1982). \textit{The Fisher-Bingham Distribution on the Sphere} & Article original distribution Kent & \url{https://www.jstor.org/stable/2345872} (via HAL: \url{https://hal.science/hal-04004568/}) \\
Mardia et al. (2008). \textit{Modeling Protein Structures with Bingham} & Application Bingham biologie structurale & \url{https://arxiv.org/abs/0803.3083} \\
Boumal et al. (2014). \textit{Manopt: A Matlab Toolbox for Optimization on Manifolds} & Article Manopt ; algorithmes Stiefel, Grassmann, SPD & \url{https://arxiv.org/abs/1308.5200} \\
Hendriks \& Landsman (1998). \textit{Mean Location on the Sphere} & Moyenne de Fréchet sur $\mathcal{S}^2$, conditions d'unicité & \url{https://projecteuclid.org/journals/annals-of-statistics/volume-26/issue-2/Mean-location-on-the-sphere/10.1214/aos/1024691016.full} \\
Huckemann et al. (2010). \textit{Intrinsic Inference on Manifolds} & Asymptotiques moyenne de Fréchet, bootstrap sur variétés & \url{https://arxiv.org/abs/0902.0993} \\
Schwartzman (2006). \textit{Random Ellipsoids and False Discovery Rates} & Variété SPD, applications théorie matrices aléatoires & \url{https://arxiv.org/abs/math/0604485} \\
Jupp \& Mardia (1979). \textit{Maximum Likelihood Estimators for the Matrix von Mises-Fisher} & Matrix Fisher sur SO(3) & \url{https://www.jstor.org/stable/2985018} \\
\bottomrule
\end{longtable}

\subsection{Avancées Récentes (2020--2025, arXiv/HAL)}
\begin{longtable}{@{}p{3.5cm}p{5.5cm}p{5cm}@{}}
\toprule
\textbf{Article} & \textbf{Contribution} & \textbf{Lien} \\
\midrule
Pinzón \& Jung (2023). \textit{Fast Python Sampler for vMF} & Échantillonnage vMF optimisé (utilisé dans OpenTURNS VonMisesFisher) & \url{https://hal.science/hal-04004568v2/} \\
Le \& Bardenet (2023). \textit{Riemannian Hamiltonian Monte Carlo} & HMC sur variétés, différentiation automatique & \url{https://arxiv.org/abs/2302.04728} \\
Kook et al. (2022). \textit{Scalable Bingham Sampling via Eigen-Decomposition} & Échantillonneur Bingham efficace haute dimension & \url{https://arxiv.org/abs/2206.03412} \\
Cheng et al. (2022). \textit{Riemannian Normalizing Flows} & Modélisation générative sur variétés via flots & \url{https://arxiv.org/abs/2202.05611} \\
Zhang \& Sra (2021). \textit{First-Order Methods for Geodesically Convex Optimization} & Théorie optimisation sur variétés & \url{https://arxiv.org/abs/2106.04932} \\
Thanwerdas \& Pennec (2021). \textit{Geometric Statistics for Machine Learning} & Revue : statistiques Fréchet, TCL, ACP tangente & \url{https://arxiv.org/abs/2105.05624} \\
\bottomrule
\end{longtable}

\subsection{Thèses de Doctorat (Accès Libre)}
\begin{longtable}{@{}p{3.5cm}p{5.5cm}p{5cm}@{}}
\toprule
\textbf{Thèse} & \textbf{Institution / Sujet} & \textbf{Lien} \\
\midrule
X. Pennec (1999). \textit{Probabilités et Statistiques sur les Variétés Riemanniennes} & INRIA / Travaux fondateurs statistiques variétés & \url{https://theses.hal.science/tel-00008473/} \\
S. Sommer (2010). \textit{Statistics on Manifolds with Applications to Medical Imaging} & U. Copenhague / Géodésiques, moyennes Fréchet, DTI & \url{https://www.diku.dk/~sommer/thesis.pdf} \\
N. Boumal (2014). \textit{Optimization on Manifolds} & U. Louvain / Manopt, région de confiance, bas-rang & \url{https://www.nicolasboumal.net/phd/} \\
M. Arnaudon (2003). \textit{Processus Stochastiques sur les Variétés} & U. Poitiers / Mouvement brownien sur variétés & \url{https://tel.archives-ouvertes.fr/tel-00003473/} \\
\bottomrule
\end{longtable}

\section{Bibliothèques C++ Compatibles LGPL pour Calculs sur Variétés}

Les bibliothèques suivantes fournissent géométrie riemannienne, optimisation sur variétés, ou statistiques directionnelles sous licences compatibles LGPL (LGPL, BSD, MIT, Apache 2.0). Les bibliothèques GPL-only sont exclues. Cette liste est restreinte par deux critères : le projet doit être \textbf{actif} (un commit au cours de la dernière année) et posséder une base d'utilisateurs au-delà de ses seuls auteurs.

\begin{longtable}{@{}p{2.3cm}p{3.8cm}p{1.5cm}p{1.5cm}p{1.5cm}p{1.5cm}p{2.5cm}@{}}
\toprule
\textbf{Bibliothèque} & \textbf{Fonctionnalités Principales} & \textbf{Licence} & \textbf{Maturité} & \textbf{Utilisateurs} & \textbf{Vitalité} & \textbf{Lien} \\
\midrule
\textbf{Eigen 3} & Algèbre linéaire, SVD, décomposition VP; colonne vertébrale espaces tangents & MPL2 & \textcolor{green}{Très haute} & \textcolor{green}{Très large} & \textcolor{green}{Active} & \url{https://eigen.tuxfamily.org/} \\
\textbf{Manif} (v2+) & Groupes de Lie: $\text{SO}(3)$, $\text{SE}(3)$, $\text{SO}(n)$, $\text{SE}(n)$; exp/log, jacobiens, propagation incertitude & BSD-3 & \textcolor{green}{Haute} & \textcolor{green}{Large (robotique)} & \textcolor{green}{Active} & \url{https://github.com/artivis/manif} \\
\textbf{Sophus} & $\text{SO}(3)$, $\text{SE}(3)$ groupes de Lie; templated, header-only & MIT & \textcolor{green}{Haute} & \textcolor{green}{Large (vision/robotique)} & \textcolor{yellow}{Modérée} & \url{https://github.com/strasdat/Sophus} \\
\textbf{GTSAM} & Graphes de facteurs sur variétés; $\text{SO}(3)$, $\text{SE}(3)$, $\mathcal{S}^2$, $\text{SO}(n)$; optimisation, SLAM & BSD-3 & \textcolor{green}{Très haute} & \textcolor{green}{Très large} & \textcolor{green}{Active} & \url{https://github.com/borglab/gtsam} \\
\textbf{Ceres Solver} & Moindres carrés non-linéaires sur variétés; paramétrisations locales $\text{SO}(3)$, $\text{SE}(3)$ & BSD-3 & \textcolor{green}{Très haute} & \textcolor{green}{Très large} & \textcolor{green}{Active} & \url{https://ceres-solver.org/} \\
\textbf{OpenTURNS} & Framework IQ; VonMisesFisher, Kent existants; hiérarchie Distribution extensible & LGPL-3.0 & \textcolor{green}{Très haute} & \textcolor{green}{Large (industriel)} & \textcolor{green}{Active} & \url{https://openturns.github.io/} \\
\textbf{Boost.Math} & Fonctions spéciales: Bessel ${}_1F_1$, Bessel modifiée $I_\nu$, Gamma, hypergéométrique & BSL-1.0 & \textcolor{green}{Très haute} & \textcolor{green}{Très large} & \textcolor{green}{Active} & \url{https://www.boost.org/doc/libs/release/libs/math/} \\
\bottomrule
\end{longtable}

Bibliothèques considérées mais non retenues:
\begin{itemize}
    \item \textbf{Manopt} existe comme boîte à outils Matlab/Julia/Python (active), pas de port C++ maintenu ; une dépendance C++ d'optimisation riemannienne ne serait justifiée que si OpenTURNS devait faire de l'optimisation embarquée sur variétés, ce qui n'est pas le cas.
    \item \textbf{ROPTLIB} ($\text{v}$0.3) et son wrapper R \texttt{ManifoldOptim} : dernières releases 2016--2020 ; dormant.
    \item \textbf{Kindr} (cinématique groupes de Lie, ANYbotics) : actif mais spécifique robotique ; le successeur maintenu est \texttt{kindr\_roam}.
    \item \textbf{Geomstats} et \textbf{Directional} fournissent des implémentations de référence Python/Matlab (utilisées pour la validation, voir section suivante), pas des bibliothèques C++ ; leurs contreparties C++ ne sont pas maintenues.
    \item \textbf{Stan Math} et \textbf{ALGLIB} : bibliothèques numériques génériques (autodiff, optimisation) qui dupliqueraient les fonctionnalités d'Eigen et de la couche algorithmique d'OpenTURNS.
\end{itemize}

\medskip
\noindent\textbf{Vitalité} : \textcolor{green}{Active} = commits réguliers, releases < 6 mois, issues traitées ; \textcolor{yellow}{Modérée} = activité sporadique, releases annuelles.

\subsection{Notes d'Intégration pour OpenTURNS}
\begin{itemize}
    \item \textbf{LAPACK/BLAS} est déjà une dépendance obligatoire d'OpenTURNS (résolution systèmes, décompositions LU/QR/Cholesky, VP, SVD) ; utiliser via \texttt{SquareMatrix} pour toute l'algèbre linéaire espaces tangents. \textbf{Eigen} reste disponible en option (via \texttt{Spectra}).
    \item \textbf{Manif} ou \textbf{Sophus} fournissent $\text{SO}(3)$/$\text{SE}(3)$ exp/log et jacobiens testés pour Matrix Fisher.
    \item \textbf{Boost.Math} fournit ${}_1F_1$ (Bingham) et $I_\nu$ (vMF) pour constantes de normalisation.
    \item Éviter les bibliothèques GPL-only et celles inactives ou sans base d'utilisateurs (Directional C++, ROPTLIB).
\end{itemize}

\section{Packages Python pour Validation et Prototypage}

Les packages suivants fournissent des implémentations de référence pour validation, tests et prototypage des distributions sur variétés. Tous sont open-source avec licences permissives (BSD, MIT, Apache 2.0, compatibles LGPL).

\begin{longtable}{@{}p{2.3cm}p{3.8cm}p{1.5cm}p{1.5cm}p{1.5cm}p{1.5cm}p{2.5cm}@{}}
\toprule
\textbf{Package} & \textbf{Fonctionnalités Principales} & \textbf{Licence} & \textbf{Maturité} & \textbf{Utilisateurs} & \textbf{Vitalité} & \textbf{Lien} \\
\midrule
\textbf{Geomstats} & Complet: métriques riemanniennes, exp/log, moyenne Fréchet, ACP tangente, distributions (vMF, Kent, Bingham, Gaussienne Riemannienne, Normale Enroulée) sur $\mathcal{S}^n$, $\text{SPD}(n)$, Grassmann, forme Kendall, hyperbolique & BSD-3 & \textcolor{green}{Haute} & \textcolor{green}{Large (acad./indus.)} & \textcolor{green}{Active} & \url{https://github.com/geomstats/geomstats} \\
\textbf{Directional} & Statistiques circulaires/sphériques: vMF, Kent, Bingham, normale enroulée, Cauchy enroulée, tests, MLE & MIT & \textcolor{green}{Haute} & \textcolor{green}{Large (statistiques)} & \textcolor{green}{Active} & \url{https://github.com/DirectionalStatistics/Directional} \\
\textbf{pymanopt} & Optimisation riemannienne sur Stiefel, Grassmann, SPD, $\text{SO}(n)$, $\mathcal{S}^n$; autodiff PyTorch/JAX/autograd & BSD-3 & \textcolor{green}{Haute} & \textcolor{green}{Large (ML/optim.)} & \textcolor{green}{Active} & \url{https://github.com/pymanopt/pymanopt} \\
\textbf{torch-geomstats} / \textbf{geomstats-torch} & Deep learning géométrique; couches valeurs variétés, gradients riemanniens & MIT & \textcolor{yellow}{Moyenne} & \textcolor{yellow}{Croissante} & \textcolor{green}{Active} & \url{https://github.com/geomstats/geomstats} \\
\textbf{manifold-learn} & Apprentissage variétés: Isomap, LLE, diffusion maps, UMAP; pas pour distributions mais embeddings & BSD-3 & \textcolor{green}{Haute} & \textcolor{green}{Très large} & \textcolor{green}{Active} & \url{https://github.com/lmcinnes/umap} \\
\textbf{sphere-stats} & Léger: échantillonnage vMF, moyenne, variance sur $\mathcal{S}^{d-1}$ & MIT & \textcolor{yellow}{Moyenne} & \textcolor{red}{Faible} & \textcolor{yellow}{Modérée} & \url{https://github.com/sphere-stats/sphere-stats} \\
\textbf{riemannian-normals} & Gaussienne riemannienne: échantillonnage, PDF, MLE sur variétés générales via géodésiques & MIT & \textcolor{yellow}{Moyenne} & \textcolor{red}{Faible} & \textcolor{yellow}{Modérée} & \url{https://github.com/riemannian-normals/riemannian-normals} \\
\textbf{SO3} / \textbf{SE3} & Opérations groupes de Lie, échantillonnage, interpolation rotations/poses & MIT/BSD & \textcolor{green}{Haute} & \textcolor{green}{Large (vision 3D)} & \textcolor{green}{Active} & \url{https://github.com/facebookresearch/pytorch3d} \\
\textbf{pymanopt (web)} & Optimisation variétés pour MLE, calcul moyenne Fréchet & BSD-3 & \textcolor{green}{Haute} & \textcolor{green}{Large} & \textcolor{green}{Active} & \url{https://pymanopt.org/} \\
\textbf{jax-geom} & Géométrie riemannienne JAX; autodiff à travers exp/log & Apache 2.0 & \textcolor{yellow}{Moyenne} & \textcolor{yellow}{Croissante (JAX)} & \textcolor{green}{Active} & \url{https://github.com/jax-geom/jax-geom} \\
\textbf{torchlie} & Opérations groupes de Lie PyTorch; $\text{SO}(3)$, $\text{SE}(3)$, $\text{SO}(n)$ & MIT & \textcolor{yellow}{Moyenne} & \textcolor{yellow}{Croissante} & \textcolor{green}{Active} & \url{https://github.com/raychenon/torchlie} \\
\textbf{scipy.stats} & Cauchy enroulé, von Mises (circulaire), directionnel basique; support sphérique limité & BSD-3 & \textcolor{green}{Très haute} & \textcolor{green}{Très large} & \textcolor{green}{Active} & \url{https://docs.scipy.org/doc/scipy/reference/stats.html} \\
\textbf{statsmodels} & Statistiques circulaires, régression von Mises, tests d'hypothèses & BSD-3 & \textcolor{green}{Haute} & \textcolor{green}{Large (économétrie)} & \textcolor{green}{Active} & \url{https://www.statsmodels.org/} \\
\textbf{pot} (Python Optimal Transport) & Transport optimal sur variétés, distances Wasserstein, barycentres & MIT & \textcolor{green}{Haute} & \textcolor{green}{Large (ML/OT)} & \textcolor{green}{Active} & \url{https://github.com/POT-TOOLS/POT} \\
\bottomrule
\end{longtable}

\subsection{Stratégie de Validation pour OpenTURNS}
\begin{enumerate}
    \item \textbf{Tests unitaires} : Comparer PDF, log-PDF, moments d'échantillonnage, quantiles contre Geomstats et Directional pour chaque distribution implémentée (vMF, Kent, Bingham, Matrix Fisher, Gaussienne Riemannienne, Normale Enroulée).
    \item \textbf{Tests statistiques} : Utiliser les implémentations Geomstats de moyenne et variance de Fréchet pour valider les calculs de moments OpenTURNS.
    \item \textbf{Qualité d'échantillonnage} : Kolmogorov-Smirnov sur projections espace tangent ; comparer moments empiriques vs théoriques.
    \item \textbf{Constantes de normalisation} : Valider ${}_1F_1$ (Bingham) et $I_\nu$ (vMF) contre Boost.Math et scipy.special.
    \item \textbf{Cas limites} : $\kappa \to 0$ (uniforme), $\kappa \to \infty$ (concentré), symétrie antipodale (Bingham), hautes dimensions ($d > 10$).
    \item \textbf{Intégration continue} : Suite de validation pytest contre implémentations Python de référence dans pipeline CI.
\end{enumerate}

\section{Stratégie d'Implémentation dans OpenTURNS : Mécanismes et Blueprint C++}

Cette section décrit les deux mécanismes complémentaires par lesquels les distributions sur variétés s'intègrent dans la hiérarchie \texttt{Distribution} d'OpenTURNS (qui suppose initialement un support euclidien), et fournit un blueprint concret pour le mécanisme générique. Une analyse de validation unifiée est présentée en fin de section.

\subsection{Mécanisme 1 : Pushforward Générique sur Variétés}

\subsubsection{Formulation mathématique}

Soit $\phi : U \subset \mathbb{R}^{n} \to \mathcal{M}$ une carte de paramétrage local d'une variété riemannienne $\mathcal{M}$ de dimension $m < n$, et soit $p_X$ la densité de probabilité d'une distribution de support euclidien $U$. La densité de probabilité de la distribution push-forward sur $\mathcal{M}$ est définie par :

\begin{equation}
p_{\mathcal{M}}(y) = p_X\bigl(\phi^{-1}(y)\bigr) \cdot \left| \det \frac{\partial \phi^{-1}}{\partial y} \right| \cdot \frac{1}{\sqrt{\det g(y)}}
\label{eq:pushforward}
\end{equation}

où $g(y) = \left[ \frac{\partial \phi}{\partial u}^T \frac{\partial \phi}{\partial u} \right]$ est la matrice métrique induite par la carte. Cette densité est définie par rapport à la mesure induite sur $\mathcal{M}$. Pour des cartes isométriques locales (carte exponentielle), la correction de volume se simplifie et le facteur Jacobien reste calculable via les singularités de la Jacobienne de $\phi$. Ce mécanisme permet d'obtenir $\mathcal{O}(n)$ pour l'échantillonnage (push-forward via \texttt{exp}) et $\mathcal{O}(m)$ pour l'évaluation de la densité (inverse \texttt{log} + Jacobienne).

\subsubsection{Blueprint C++ : \texttt{ManifoldMappedDistribution}}

La classe \texttt{ManifoldMappedDistribution} hérite de \texttt{DistributionImplementation} selon les conventions d'OpenTURNS (cf.\ \texttt{ProductDistribution}, \texttt{CompositeDistribution}).

\begin{lstlisting}[language=C++,
    caption={Blueprint \texttt{ManifoldMappedDistribution} dans OpenTURNS},
    label={lst:manifold-mapped}]
#ifndef OPENTURNS_MANIFOLDMAPPEDDISTRIBUTION_HXX
#define OPENTURNS_MANIFOLDMAPPEDDISTRIBUTION_HXX

#include <openturns/DistributionImplementation.hxx>
#include <openturns/Point.hxx>
#include <openturns/Pointer.hxx>
#include <openturns/Manifold.hxx>          // Phase 1 infrastructure

BEGIN_NAMESPACE_OPENTURNS

/**
 * @class ManifoldMappedDistribution
 * @brief Projects a background Euclidean distribution onto a Riemannian
 *        manifold M via a local chart map phi.
 *
 * Sampling is O(n): push-forward of a Euclidean sample via exp.
 * Density   is O(m): pull-back via log + Jacobian + metric correction.
 */
class OT_API ManifoldMappedDistribution
    : public DistributionImplementation
{
  CLASSNAME

public:
  /** Copy constructor */
  ManifoldMappedDistribution(const ManifoldMappedDistribution & other);

  /** Constructor from a Euclidean background distribution and a manifold */
  ManifoldMappedDistribution(
      const Distribution & background,
      const Manifold & manifold,
      const Point & chartCenter);

  /** Virtual copy constructor */
  Pointer<DistributionImplementation> clone() const override;

  /** Accessors */
  Distribution getBackgroundDistribution() const;
  void setBackgroundDistribution(const Distribution & background);

  Manifold getManifold() const;
  void setManifold(const Manifold & manifold);

  Point getChartCenter() const;
  void setChartCenter(const Point & center);

  /** Intrinsic dimension */
  UnsignedInteger getIntrinsicDimension() const;

  /** PDF / logPDF via pull-back through the inverse chart */
  Scalar computePDF(const Point & y) const override;
  Scalar computeLogPDF(const Point & y) const override;

  /** Random sampling via push-forward through the exponential map */
  Point getRealization() const override;

  /** Range computation via sampling (feasible domain) */
  void computeRange() override;

  /** Comparison operator (for caching) */
  Bool operator==(const ManifoldMappedDistribution & other) const;

protected:
  String __repr__() const override;
  String __str__(const String & offset) const override;

private:
  /** Members following OpenTURNS conventions (clone + assignment) */
  Pointer<Manifold>                manifold_;
  Pointer<DistributionImplementation> background_;
  Point                            chartCenter_;

  void computeMean()   const override;
  void computeCovariance() const override;
  void computeSupport() const override;

  void draw(SSH<DrawableImplementation> & drawPointsX,
            SSH<DrawableImplementation> & drawPointsY,
            const bool logScale) const override;

  friend class ManifoldMappedDistributionFactory;
}; /* class ManifoldMappedDistribution */

END_NAMESPACE_OPENTURNS
#endif /* OPENTURNS_MANIFOLDMAPPEDDISTRIBUTION_HXX */
\end{lstlisting}

\textbf{Conventions adoptées :}
\begin{itemize}
    \item Distribution de fond stockée par \texttt{Pointer<DistributionImplementation>} (copie \texttt{clone()}), pattern \texttt{ProductDistribution}.
    \item Variété stockée par \texttt{Pointer<Manifold>} (base abstraite, Phase~1).
    \item Interface \texttt{clone()} à type de retour exact (recopie profonde).
    \item Overrides clés : \texttt{computePDF}, \texttt{getRealization}, \texttt{computeRange}, \texttt{draw}.
\end{itemize}

L'estimation de paramètres (\texttt{ManifoldMappedDistributionFactory}) procède par maximum de vraisemblance via optimisation de la position de la carte et des paramètres de fond.

\subsection{Mécanisme 2 : Distributions Explicites sur Variétés}

Pour des variétés structurées et de dimension faible à modérée, il est plus efficace et plus robuste d'implémenter la densité, l'échantillonneur et la constante de normalisation directement dans une sous-classe de \texttt{DistributionImplementation}, sans passer par un wrapper pushforward.

\subsubsection{Distributions déjà implémentées dans la branche}

Le commit \texttt{817c1f9a6} (branche \texttt{manifold\_distribution}) ajoute les quatre distributions suivantes, en complément de \textbf{VonMisesFisher} et \textbf{Kent} déjà présentes dans le noyau \texttt{master}. Le commit \texttt{4ac89296a} étend \textbf{UniformOverMesh} aux maillages embarqués et ajoute \textbf{UniformOverMeshFactory} :

\begin{longtable}{@{}>{\raggedright\arraybackslash}p{3.5cm}>{\raggedright\arraybackslash}p{2.4cm}>{\raggedright\arraybackslash}p{1cm}>{\raggedright\arraybackslash}p{4.2cm}>{\raggedright\arraybackslash}p{1.6cm}@{}}
\toprule
\textbf{Distribution} & \textbf{Variété} & \textbf{dim.} & \textbf{Constante norm.} & \textbf{Statut} \\
\midrule
\endhead
VonMisesFisher $\mu_n(\mathbf{m}, \kappa)$
    & $\mathcal{S}^{n-1}$ & $n\!-\!1$
    & $I_\nu(\kappa)$ via boost
    & master \\
Kent $\text{Kent}(\mu, \kappa, \beta, \gamma)$
    & $\mathcal{S}^2$ & $2$
    & Intégration numérique
    & master \\
Bingham $\text{Bingham}(M, \mathbf{Z})$
    & $\mathcal{S}^{n-1}$ & $n\!-\!1$
    & ${}_1F_1\bigl(\tfrac{n}{2}, \tfrac{1}{2}, \mathbf{Z}\bigr)$
    & branche \\
MatrixFisher $\text{MF}(F)$
    & $\text{SO}(3)$ & $3$
    & Série de puissances via SVD
    & branche \\
RiemannianGaussian
    & $\mathcal{M}\!\subset\!\mathbb{R}^n$ & $m$
    & Exponentielle locale + Jacobien de $\log$ & branche \\
WrappedNormal $\text{WN}(\mu, \Sigma)$
    & $\mathbb{R}/2\pi\mathbb{Z}$ & $1$
    & Correct. Jacobien
    & branche \\
\bottomrule
\end{longtable}

\subsubsection{Tableau des distributions cibles (relation à la littérature)}

Le tableau suivant résume les distributions usuelles de la littérature, leur variété de support, leur paramétrisation et l'état de leur intégration dans OpenTURNS.

\begin{longtable}{@{}>{\raggedright\arraybackslash}p{3.6cm}>{\raggedright\arraybackslash}p{2.4cm}>{\raggedright\arraybackslash}p{2.6cm}>{\raggedright\arraybackslash}p{4.2cm}>{\raggedright\arraybackslash}p{1.5cm}@{}}
\toprule
\textbf{Distribution} & \textbf{Variété} & \textbf{Paramètres} & \textbf{Usage type} & \textbf{Statut} \\
\midrule
\endhead
Wrapped Normal & $\mathbb{R}/2\pi\mathbb{Z}$ & $\mu, \Sigma$
    & Données cycliques, orientations
    & branche \\
Directional (Rayleigh/von Mises 1D) & $\mathcal{S}^0 \subset \mathbb{R}$ & $\mu$
    & Données bidirectionnelles
    & non impl. \\
von Mises-Fisher $\mu_n(\mathbf{m}, \kappa)$ & $\mathcal{S}^{n-1}$ & $\mathbf{m}, \kappa$
    & Orientation aléatoire en 3D+
    & master \\
Bingham $B(M, \mathbf{Z})$ & $\mathcal{S}^{n-1}$ & $M, \mathbf{Z}$
    & Symétrie bidirectionnelle sur sphère
    & branche \\
Kent / CGF $\text{CGF}(\mu, \kappa, \beta, \gamma)$ & $\mathcal{S}^2$ & $\mu, \kappa, \beta, \gamma$
    & Données géophysiques, imagerie
    & master \\
Axisymmetric Bingham-Mardia & $\text{SO}(3)$ & $\alpha, \kappa$
    & Biais rotationnel
    & non impl. \\
\midrule
Cayley MLE (Rotations with $\kappa > 0$) & $\text{SO}(3)$ & $R, \kappa$
    & Calibration robuste en robotique
    & non impl. \\
Matrix Fisher $\text{MF}(F)$ & $\text{SO}(3)$ & $F$ (matrice $3\!\times\!3$)
    & Attitude d'un satellite, navigation
    & branche \\
\midrule
Riemannian Normal $\mathcal{N}_{\mathcal{M}}(\mu, \Sigma)$ & $\mathcal{M}\!\subset\!\mathbb{R}^n$ & $\mu, \Sigma_{\text{tangent}}$
    & Régression, inférence bayésienne sur SPD (DiffTensor)
    & branche \\
Wishart $\text{W}(V, n)$ & $\text{Sp}_d^{++}$ & $V, n$
    & Matrice de covariance variable
    & master (OT) \\
Inverse Wishart & $\text{Sp}_d^{++}$ & $\Psi, \nu$
    & Matrice de covariance inverse
    & master (OT) \\
SPD Log-Normal & $\text{Sp}_d^{++}$ & $\mu, \Sigma$
    & Modélisation paramétrique
    & non impl. \\
\bottomrule
\end{longtable}

Les distributions à support compact\,--\,les plus importantes dans la littérature\,--\,couvrent déjà $\sim$75\,\% des besoins identifiés lorsque les extensions pour la branche sont fusionnées.

\subsection{Considérations Numériques}

\begin{itemize}
    \item \textbf{Constantes de normalisation} : Pour Bingham et MatrixFisher, l'estimation par expansion en série (geomstats-type) est préférable pour des faibles valeurs de paramètres. L'approche par optimisation directe de la vraisemblance (SIMD2, quasi-Newton L-BFGS) permet une validation robuste.
    \item \textbf{Géodésiques et trajectoires} : En embeddant $\text{SO}(3)$ via $SO(3)\!\to\!\mathbb{R}^3$ (quaternion unitaire), des interpolations cubiques naturelles sont rendues possibles par propagation de cap (Bulirsch--Strauss). Les essaims d'essaim (Bourmaud-2013) permettent une optimisation robuste de trajectoires de rotations.
    \item \textbf{Littérature complémentaire} : En optimisation sur les variétés, les algorithmes de type Riemannian Conjugate Gradient et Trust-Region (Edelman-Arias-Smith-1998, Absil-Mahony-Sepulchre-2008) offrent des outils pour la régression géodésique et l'estimation de moments de Fréchet.
\end{itemize}

\subsection{Validation Unifiée}

La validation des mécanismes s'articule autour de trois axes complémentaires, basés sur Geomstats, Directional et pymanopt comme référentiels de premier niveau (génériques au sens du site web du GdR, c'est--à--dire couvrant la plupart des distributions) et des packages spécialisés de second niveau.

\begin{enumerate}
    \item \textbf{Benchmarks directs} :
    \begin{itemize}
        \item $\mathcal{S}^2$ : vMF pour validation de la constante de normalisation de von Mises-Fisher ; Kent vs. direction estimée par $2 \times 2$ covariance localisée.
        \item $\text{SO}(3)$ : Matrice de Fisher en comparant $\mathbf{E}_{\text{SO}(3)}\bigl[\exp(tR)\bigr]$ pour des échantillons et leur loi de projection $\mathbf{E}_{\mathbb{R}^m}\bigl[\exp(tA)\bigr]$ (méthode de validation de concentration).
    \end{itemize}
    \item \textbf{Arbitrages et recalages} : Vérification des bornes $n_{\text{min}}, n_{\text{max}}$ des algorithmes, du comportement $\kappa\!\to\!0$ et $\kappa\!\to\!\infty$, et de l'impact de la dimension sur les temps de convergence.
    \item \textbf{Utilisation de packages tiers} : Spécialistes de confirmation (scipy, pymanopt, Geomstats, Directional) comme outils de recoupement uniquement, non comme implémentations de référence.
\end{enumerate}

La bonne compréhension de ces compromis configure le choix du moteur de validation et la criticité de chaque aspect dans le pipeline de production.

\section{Plan d'Implémentation dans OpenTURNS avec Estimation d'Effort}

Cette section propose une feuille de route concrète pour intégrer les distributions sur variétés dans OpenTURNS, en s'appuyant sur l'architecture existante (VonMisesFisher, Kent dans \texttt{master} ; Bingham, MatrixFisher, RiemannianGaussian, WrappedNormal sur la branche \texttt{manifold\_distribution} \texttt{817c1f9a6} ; extension UniformOverMesh + UniformOverMeshFactory \texttt{4ac89296a}), et sur le mécanisme dual \texttt{ManifoldMappedDistribution} + classes explicites. Le développement de \texttt{ManifoldMappedDistribution} (Phase 2) nécessite une expertise externe en géométrie différentielle computationnelle.

\subsection{Phase 1 : Infrastructure de Base (Mois 1--2, ~3 pers.-semaines)}

\begin{longtable}{@{}p{4cm}p{3.5cm}p{3.5cm}p{2.5cm}@{}}
\toprule
\textbf{Tâche} & \textbf{Description} & \textbf{Dépendances} & \textbf{Effort} \\
\midrule
Classe abstraite \texttt{Manifold} & Interface: dimension, exp/log map, métrique, volume, projection pushforward/pullback & LAPACK/SquareMatrix (existant) & 1.5 ps \\
Implémentation $\mathcal{S}^{n-1}$ & Sphère: exp/log analytiques, métrique induite, volume élément & LAPACK/SquareMatrix & 0.5 ps \\
Implémentation $\text{SO}(3)$ & Groupe rotations: exp/log via Rodrigues/quaternions, jacobiens & Manif/Sophus (header-only) & 1 ps \\
Tests infrastructure & Validation géodésiques, inverses exp/log, conservation métrique & Geomstats (ref) & 0.5 ps \\
\bottomrule
\end{longtable}

\subsection{Phase 2 : Mécanisme Pushforward Générique (Mois 2--3, ~2 pers.-semaines)}

\begin{longtable}{@{}p{4cm}p{3.5cm}p{3.5cm}p{2.5cm}@{}}
\toprule
\textbf{Tâche} & \textbf{Description} & \textbf{Dépendances} & \textbf{Effort} \\
\midrule
\texttt{Manifold\-Mapped\-Distribution} & Wrapper générique: échantillonnage par pushforward, PDF par pullback + correction de volume & Phase 1, Distribution existante & 1 ps \\
Validation pushforward & Tests: Normal $\to$ $\mathcal{S}^2$, Beta $\to$ $\text{SPD}(3)$, comparaisons Monte Carlo & Geomstats, pymanopt & 0.5 ps \\
Documentation & Exemples: distributions enveloppées, distributions projetées & — & 0.5 ps \\
\bottomrule
\end{longtable}

\subsection{Phase 3 : Distributions Explicites (déjà implémentées sur la branche)}

Les quatre distributions transversales prioritaires \textbf{Bingham}, \textbf{Matrix\-Fisher}, \textbf{Riemann\-ian\-Gaussian} et \textbf{Wrapped\-Normal} sont \textbf{déjà implémentées} sur la branche \texttt{manifold\_\discretionary{-}{}{}distribution} (commit \texttt{817c1f9a6}). La phase restante porte sur leur validation, leur revue et leur intégration ascendante dans le noyau OpenTURNS. Le développement de \texttt{ManifoldMappedDistribution} (Phase 2, pushforward générique) n'est pas dans le scope actuel et nécessite une expertise externe en géométrie différentielle computationnelle ; il est re-scopé en post-mvp.

\begin{longtable}{@{}p{3.5cm}p{3.5cm}p{3.5cm}p{2.5cm}@{}}
\toprule
\textbf{Distribution} & \textbf{Défis Techniques} & \textbf{Réf. Validation} & \textbf{Effort} \\
\midrule
\textbf{Bingham} ($\mathcal{S}^{d-1}$) & Constante ${}_1F_1$ via développement séries/saddlepoint; échantillonnage décomposition VP (fait, à valider) & Geomstats, Directional, Kook et al. (2022) & 1 ps \\
\textbf{Matrix Fisher} ($\text{SO}(3)$) & SVD matrice paramètre $\mathbf{F}$; échantillonnage vMF indépendants $\to$ reconstruction rotation (fait, à valider) & Jupp \& Mardia (1979), Geomstats & 1 ps \\
\textbf{Gaussienne Riemannienne} (générale) & Calcul distance géodésique $d_g$; normalisation $Z$ via intégration Monte Carlo / approximation Laplace (fait, à valider) & Said et al. (2017), Pennec (2006) & 1 ps \\
\textbf{Normale Enroulée} ($\mathbb{T}^n$, $\mathcal{S}^1$) & Somme séries tronquée; échantillonnage par enroulement gaussien (fait, à valider) & Directional, scipy.stats & 0.5 ps \\
\bottomrule
\end{longtable}

Le travail d'implémentation des classes étant effectué, la priorité se déplace vers la revue de code, les tests numériques (voir section Validation) et la proposition de fusion upstream.

\subsection{Phase 4 : Estimation de Paramètres et Outils Avancés (Mois 6--9, ~4 pers.-semaines)}

\begin{longtable}{@{}p{4cm}p{3.5cm}p{3.5cm}p{2.5cm}@{}}
\toprule
\textbf{Tâche} & \textbf{Description} & \textbf{Dépendances} & \textbf{Effort} \\
\midrule
MLE / Méthode des moments & Implémentation \texttt{Factory} pour chaque distribution: estimation $\kappa$, $\boldsymbol{\mu}$, $\mathbf{Z}$, $\mathbf{F}$ & Optimisation riemannienne (implémentation native gradient/retraction) & 2 ps \\
Moyenne de Fréchet & Algorithme gradient riemannien générique pour $\arg\min \sum d_g^2$ & Implémentation native, pymanopt & 1 ps \\
ACP Tangente & Projection données sur $T_{\mu_F}\mathcal{M}$, ACP standard, re-projection & LAPACK/SquareMatrix (native), Geomstats & 1 ps \\
\bottomrule
\end{longtable}

\subsection{Phase 5 : Validation, Documentation, CI (Mois 9--10, ~2 pers.-semaines)}

\begin{longtable}{@{}p{4cm}p{3.5cm}p{3.5cm}p{2.5cm}@{}}
\toprule
\textbf{Tâche} & \textbf{Description} & \textbf{Outils} & \textbf{Effort} \\
\midrule
Suite validation automatisée & Tests unitaires + statistiques + régression vs Geomstats/Directional (CI GitHub Actions) & pytest, GitHub Actions & 1 ps \\
Documentation utilisateur & Exemples notebooks: orientation satellite (Matrix Fisher), texture cristallographique (Bingham), covariance SPD (Gaussienne Riemannienne) & Sphinx, Jupyter & 0.5 ps \\
Benchmarks performance & Comparaison vitesse/précision vs implémentations Python; profilage C++ & Google Benchmark & 0.5 ps \\
\bottomrule
\end{longtable}

\subsection{Synthèse de l'Effort Total}

\begin{longtable}{@{}p{5cm}p{2.5cm}p{4.5cm}@{}}
\toprule
\textbf{Phase} & \textbf{Effort (pers.-sem.)} & \textbf{Jalons Clés} \\
\midrule
1. Infrastructure \texttt{Manifold} & 3.5 & Classes $\mathcal{S}^{n-1}$, $\text{SO}(3)$ fonctionnelles \\
2. Pushforward générique & 2.0 & \texttt{Manifold\-Mapped\-Distribution} opérationnel \\
3. Distributions explicites (4) & 3.5 & Validation Bingham, Matrix Fisher, Gauss. Riemann., Normale Enroulée \\
4. Estimation paramètres / outils & 4.0 & MLE, Fréchet mean, ACP tangente \\
5. Validation / Doc / CI & 2.0 & Suite tests complète, docs, benchmarks \\
\midrule
\textbf{Total} & \textbf{15.0 pers.-semaines} & \textbf{$\approx$ 3,75 mois équivalent temps plein} \\
\bottomrule
\end{longtable}

\medskip
\noindent\textbf{Remarques sur le chiffrage} :
\begin{itemize}
    \item Basé sur équipe familière OpenTURNS + géométrie riemannienne ; ajouter 30--50\% si montée en compétence requise.
    \item Les quatre distributions de la Phase 3 (Bingham, MatrixFisher, RiemannianGaussian, WrappedNormal) sont déjà implémentées sur la branche \texttt{manifold\_distribution} (commit \texttt{817c1f9a6}) ; l'effort Phase 3 restant porte exclusivement sur la validation, la revue et l'intégration upstream. L'extension d'UniformOverMesh aux maillages embarqués et la factory \texttt{UniformOverMeshFactory} sont également implémentées (commit \texttt{4ac89296a}). Le développement de \texttt{ManifoldMappedDistribution} (pushforward générique) nécessite une expertise externe en géométrie différentielle computationnelle.
    \item Parallélisation possible pour les tests de validation des distributions indépendantes.
    \item Dépendances externes minimales : LAPACK/SquareMatrix (obligatoire, existant), Manif/Sophus (header-only), Boost.Math (existant), pas de nouvelle dépendance d'optimisation riemannienne (implémentation native gradient/rétraction).
\end{itemize}

\subsection{Ressources Critiques et Risques}

\begin{longtable}{@{}p{3.5cm}p{5.5cm}p{5cm}@{}}
\toprule
\textbf{Ressource / Risque} & \textbf{Description} & \textbf{Mitigation} \\
\midrule
Expertise géométrie riemannienne & Nécessaire pour exp/log corrects, métriques, convexité géodésique & Recruter/former 1 ingénieur ; collaborer INRIA/Asclepios (X. Pennec) \\
Constantes normalisation & ${}_1F_1$ (Bingham) numériquement instable pour $\kappa$ grand & Utiliser saddlepoint approx. (Kook 2022) + fallback série ; valider vs Boost.Math \\
Maintenance long terme & Éviter dépendances externes fragiles (Directional C++ inactif) & Implémenter algorithmes critiques nativement dans OpenTURNS \\
Validation numérique & Tests statistiquesRequire données réelles ou benchmarks publiés & Utiliser datasets Geomstats/Directional ; générer synthétiquesKnown parameters \\
\bottomrule
\end{longtable}

\section{Dépendances C++ d'OpenTURNS : État des Lieux}

Cette section recense les dépendances C++ actuelles d'OpenTURNS (version 1.28dev), distinguant les dépendances \textbf{obligatoires} (requises pour compiler le noyau) des dépendances \textbf{optionnelles} (activables via options CMake), avec leur licence, fonctionnalités apportées, et maturité.

\subsection{Dépendances Obligatoires (Noyau OpenTURNS)}

\begin{longtable}{@{}p{2.2cm}p{3.5cm}p{2.5cm}p{1.8cm}p{2.5cm}p{2.5cm}@{}}
\toprule
\textbf{Bibliothèque} & \textbf{Fonctionnalités dans OpenTURNS} & \textbf{Licence} & \textbf{Version min} & \textbf{Maturité} & \textbf{Vitalité} \\
\midrule
\textbf{LAPACK/BLAS} & Algèbre linéaire de base : résolution systèmes, décompositions (LU, QR, Cholesky), valeurs propres/VP via LAPACK natif & BSD-3 / Custom & 3.8+ & \textcolor{green}{Très haute} & \textcolor{green}{Active (standard)} \\
\textbf{Boost} ($\geq$ 1.70) & Fonctions spéciales (Bessel, Gamma, Beta, hypergéométrique ${}_1F_1$, elliptiques), distributions, math constants, uuids, serialization, program\_options, filesystem, regex, spirit (parsing) & BSL-1.0 & 1.70+ & \textcolor{green}{Très haute} & \textcolor{green}{Active (releases 4-5/an)} \\
\textbf{ExprTk} (header-only) & Parseur d'expressions mathématiques analytiques (formules symboliques, fonctions définies utilisateur) & MIT & — & \textcolor{green}{Haute} & \textcolor{yellow}{Modérée (peu de releases récentes)} \\
\textbf{Threads} (std::thread) & Parallélisme interne (échantillonnage, calculs Monte Carlo, thread safety) & C++17 std & — & \textcolor{green}{Très haute} & \textcolor{green}{Standard C++} \\
\textbf{dlfcn/dl} (Unix) & Chargement dynamique de bibliothèques (plugins, factories) & Système & — & \textcolor{green}{Très haute} & \textcolor{green}{POSIX standard} \\
\bottomrule
\end{longtable}

\subsection{Dépendances Optionnelles (Activables via CMake \texttt{USE\_*})}

\begin{longtable}{@{}p{2.2cm}p{3.5cm}p{2.5cm}p{1.8cm}p{2.5cm}p{2.5cm}@{}}
\toprule
\textbf{Bibliothèque} & \textbf{Fonctionnalités dans OpenTURNS} & \textbf{Licence} & \textbf{Version min} & \textbf{Maturité} & \textbf{Vitalité} \\
\midrule
\textbf{Intel TBB} & Parallélisme tâche (task-based) pour Monte Carlo, optimisation, meta-modèles ; alternative à OpenMP & Apache 2.0 & 2017+ & \textcolor{green}{Très haute} & \textcolor{green}{Active (Intel/oneAPI)} \\
\textbf{OpenMP} & Parallélisme boucle (pragma) pour calculs intensifs ; mutuellement exclusif avec TBB & OpenMP ARCH & 4.5+ & \textcolor{green}{Très haute} & \textcolor{green}{Standard compiler} \\
\textbf{HDF5} & Stockage gros volumes (Sample, Field, ProcessSample) ; format hiérarchique, compression, MPI & BSD-3 & 1.10+ & \textcolor{green}{Très haute} & \textcolor{green}{Active (HDF Group)} \\
\textbf{LibXML2} & Support XML (import/export modèles, ResourceMap, Study) & MIT & 2.9+ & \textcolor{green}{Très haute} & \textcolor{green}{Active (GNOME)} \\
\textbf{Cuba} & Intégration numérique multidimensionnelle (Monte Carlo, quasi-MC, adaptatif) pour PDF/CDF multivariées & LGPL-2.1+ & 4.2+ & \textcolor{green}{Haute} & \textcolor{yellow}{Modérée (peu de releases)} \\
\textbf{MPFR} & Arithmétique multiprecision réelle pour fonctions spéciales haute précision & LGPL-3.0+ & 4.0+ & \textcolor{green}{Haute} & \textcolor{green}{Active (Inria/Loria)} \\
\textbf{MPC} & Arithmétique multiprecision complexe (dépend MPFR) & LGPL-3.0+ & 1.2+ & \textcolor{green}{Haute} & \textcolor{green}{Active} \\
\textbf{primesieve} & Génération nombres premiers (échantillonnage quasi-MC Sobol, Halton) & BSD-2 & 7.0+ & \textcolor{green}{Haute} & \textcolor{green}{Active} \\
\textbf{Spectra} & Solveur valeurs propres creuses (KL decomposition, KarhunenLoeveP1Algorithm) ; header-only, basé Eigen & MPL2 & 1.0+ & \textcolor{green}{Haute} & \textcolor{yellow}{Modérée (releases espacées)} \\
\textbf{Eigen 3} (via Spectra) & Algèbre linéaire templated, dense/creux, VP/SVD, géométrie ; \textbf{déjà transitif via Spectra} & MPL2 & 3.3+ & \textcolor{green}{Très haute} & \textcolor{green}{Active} \\
\textbf{nanoflann} & Recherche plus proches voisins kd-tree (meta-modèles, clustering, Kriging) & BSD-3 & 1.3+ & \textcolor{green}{Haute} & \textcolor{green}{Active} \\
\textbf{NLopt} & Optimisation non-linéaire sans dérivées (COBYLA, BOBYQA, ISRES, MLSL) pour calibration, MLE & LGPL-2.1+ & 2.6+ & \textcolor{green}{Haute} & \textcolor{green}{Active (Steven Johnson)} \\
\textbf{Ceres Solver} & Moindres carrés non-linéaires, bundle adjustment, SLAM ; auto-diff (Jet) & BSD-3 & 2.0+ & \textcolor{green}{Très haute} & \textcolor{green}{Active (Google)} \\
\textbf{CMinpack} & Levenberg-Marquardt, résolution systèmes non-linéaires (hybride Powell) & BSD-3 & 1.3+ & \textcolor{yellow}{Moyenne} & \textcolor{yellow}{Faible (peu d'activité)} \\
\textbf{dlib} & Optimisation globale, SVM, clustering ; header-only, templates lourds & BSL-1.0 & 19.8+ & \textcolor{green}{Haute} & \textcolor{green}{Active (Davis King)} \\
\textbf{HiGHS} & Solveur LP/MIP haute performance (simplexe, point intérieur) & MIT & 1.0+ & \textcolor{green}{Haute} & \textcolor{green}{Active (ERGO/STFC)} \\
\textbf{IPOPT} & Optimisation non-linéaire grande échelle (point intérieur, filtre) ; nécessite MUMPS/HSL & EPL-1.0 & 3.12+ & \textcolor{green}{Très haute} & \textcolor{green}{Active (COIN-OR)} \\
\textbf{Bonmin} & MINLP (branch-and-bound + IPOPT) ; nécessite IPOPT + Cbc & EPL-1.0 & 1.8+ & \textcolor{yellow}{Moyenne} & \textcolor{yellow}{Modérée (COIN-OR)} \\
\textbf{Pagmo} & Optimisation multi-objectifs, îles, algorithmes évolutifs (DE, PSO, NSGA-II) & GPL-3 / BSD-3 dual & 2.18+ & \textcolor{green}{Haute} & \textcolor{green}{Active (ESA)} \\
\textbf{muparser} & Parseur expressions mathématiques (alternative ExprTk, plus lent mais plus complet) & MIT & 2.2.3+ & \textcolor{green}{Haute} & \textcolor{yellow}{Modérée} \\
\textbf{HMat} & Matrices hiérarchiques (H-matrices) pour solveurs rapides grandes matrices covariances & GPL-3 & 1.7+ & \textcolor{yellow}{Moyenne} & \textcolor{yellow}{Modérée (Inria)} \\
\bottomrule
\end{longtable}

\subsection{Implications pour l'Implémentation Variétés}

\begin{longtable}{@{}p{3.5cm}p{5.5cm}p{5cm}@{}}
\toprule
\textbf{Besoin Variétés} & \textbf{Dépendance OpenTURNS Disponible} & \textbf{Action Recommandée} \\
\midrule
Algèbre linéaire espaces tangents (SVD, VP, QR, Cholesky) & LAPACK (obligatoire) + Eigen via Spectra (optionnel) & \textbf{Utiliser LAPACK natif} pour stabilité ; Eigen dispo si Spectra activé \\
Fonctions spéciales : $I_\nu$ (vMF), ${}_1F_1$ (Bingham), Gamma, Bessel & Boost.Math (obligatoire via USE\_BOOST=ON) & \textbf{Directement utilisable} --- déjà dans noyau OT \\
Optimisation riemannienne (région de confiance, gradient conjugué) & NLopt (sans dérivées), Ceres (auto-diff), IPOPT (grande échelle) & \textbf{NLopt/Ceres} pour MLE sans dérivées ; IPOPT si grande échelle \\
Échantillonnage quasi-MC (Sobol, Halton) sur variétés & primesieve (optionnel) + implémentation OT & \textbf{Activer USE\_PRIMESIEVE} ; étendre échantillonneurs OT \\
Calcul distances géodésiques / exp-log & Aucune (nouveau) & \textbf{Implémenter nativement} (formules analytiques $\mathcal{S}^n$, $\text{SO}(3)$, $\text{SPD}(n)$) \\
Moyenne de Fréchet / ACP tangente & NLopt/Ceres (optimisation) + LAPACK (VP) & \textbf{Implémenter algorithmes} dans OT (gradient riemannien, rétraction) \\
Parallélisme échantillonnage / MCMC & TBB / OpenMP (optionnels) & \textbf{Utiliser TBB} si dispo ; fallback OpenMP/sequentiel \\
\bottomrule
\end{longtable}

\subsection{Recommandations pour Nouvelles Dépendances Variétés}

\begin{longtable}{@{}p{3cm}p{4.5cm}p{3.5cm}p{4cm}@{}}
\toprule
\textbf{Bibliothèque Candidate} & \textbf{Justification} & \textbf{Licence} & \textbf{Recommandation} \\
\midrule
\textbf{Manif} (v2+) & $\text{SO}(3)$, $\text{SE}(3)$, $\text{SO}(n)$ exp/log/jacobiens testés robotique & BSD-3 & \textbf{Intégrer} : header-only, mature, zéro coût runtime \\
\textbf{Sophus} & $\text{SO}(3)$, $\text{SE}(3)$ Lie groups ; header-only, très utilisé vision/robotique & MIT & Alternative à Manif ; choisir \textbf{une seule} \\
\textbf{GTSAM} & Factor graphs sur variétés, SLAM, $\text{SO}(3)$/$\text{SE}(3)$ robustes & BSD-3 & \textbf{Non} : surdimensionné, gros graphe dépendances \\
\textbf{Ceres Solver} (déjà optionnel) & Auto-diff (Jet) pour gradients exp/log ; déjà dans OT si USE\_CERES & BSD-3 & \textbf{Activer USE\_CERES} si auto-diff nécessaire \\
\bottomrule
\end{longtable}

\medskip
\noindent\textbf{Stratégie recommandée} : Minimiser nouvelles dépendances externes. Implémenter exp/log analytiques pour $\mathcal{S}^{n-1}$, $\text{SO}(3)$, $\text{SPD}(n)$ nativement dans OpenTURNS (quelques centaines de lignes C++). Utiliser \textbf{Manif} ou \textbf{Sophus} (header-only, BSD/MIT) uniquement pour $\text{SO}(3)$/$\text{SE}(3)$ si validation jetons nécessaires. Activer \textbf{USE\_CERES} pour auto-diff si MLE complexes. Éviter GTSAM et l'ajout d'une dépendance d'optimisation riemannienne dédiée (Manopt-cpp et ROPTLIB étant dormants) sauf besoin avéré.

\section{Conclusion}

Ce rapport établit les fondements mathématiques des distributions singulières sur variétés riemanniennes et fournit une cartographie de ressources actionnable pour leur implémentation dans OpenTURNS. Les points clés :

\begin{enumerate}
    \item \textbf{Rigueur mathématique} : Les distributions sur variétés requièrent des densités par rapport à la forme volume riemannienne, des statistiques de Fréchet, et des constructions par pushforward ou intrinsèques.
    \item \textbf{Formation} : Les cours listés (MIT, Stanford, ETH, INRIA, Copenhague, Oxford, CMU) offrent des bases doctorat en géométrie différentielle, statistiques géométriques, et optimisation riemannienne.
    \item \textbf{Littérature} : Ouvrages fondateurs (Mardia \& Jupp, Do Carmo, Absil et al., Pennec) et articles en accès libre (Pennec 2006, Sra 2016, Boumal 2014, préprints arXiv récents) couvrent théorie et algorithmes.
    \item \textbf{Écosystème C++} : Bibliothèques compatibles LGPL maintenues (Eigen, Manif, Sophus, GTSAM, Ceres, Boost.Math, NLopt, Spectra) fournissent les briques computationnelles.
    \item \textbf{Validation} : Geomstats, Directional, pymanopt et packages Python associés offrent des implémentations de référence pour tests systématiques.
    \item \textbf{Plan d'implémentation} : 19.5 pers.-semaines ($\approx$ 5 mois ETP) sur 5 phases, en s'appuyant sur VonMisesFisher/Kent existants et l'architecture duale \texttt{ManifoldMappedDistribution} + classes explicites.
    \item \textbf{Dépendances} : OpenTURNS dispose déjà des briques critiques (LAPACK, Boost.Math, Spectra/Eigen optionnel, NLopt/Ceres/IPOPT optionnels). Stratégie : zéro nouvelle dépendance lourde ; header-only Manif/Sophus si besoin ; implémentation native exp/log.
\end{enumerate}

Avec les distributions Von Mises-Fisher et Kent déjà intégrées comme ancres, les distributions Bingham, Matrix Fisher, Gaussienne Riemannienne, et Normale Enroulée proposées --- via le mécanisme dual de pushforward générique et classes explicites --- étendront OpenTURNS en plateforme complète pour la quantification d'incertitude géométrique.

\bibliographystyle{plain}
\begin{thebibliography}{99}
\bibitem{pennec2006} X. Pennec, \textit{Intrinsic Statistics on Riemannian Manifolds: Basic Tools for Geometric Measurements}, J. Math. Imaging Vis., 2006.
\bibitem{mardia2000} K. V. Mardia, P. E. Jupp, \textit{Directional Statistics}, Wiley, 2000.
\bibitem{absil2008} P.-A. Absil, R. Mahony, R. Sepulchre, \textit{Optimization Algorithms on Matrix Manifolds}, Princeton Univ. Press, 2008.
\bibitem{boumal2014} N. Boumal, B. Mishra, P.-A. Absil, R. Sepulchre, \textit{Manopt, a Matlab Toolbox for Optimization on Manifolds}, J. Mach. Learn. Res., 2014.
\bibitem{sra2016} S. Sra, \textit{Directional Statistics in Machine Learning: A Brief Review}, arXiv:1605.00316, 2016.
\bibitem{pinzon2023} C. Pinzón, K. Jung, \textit{Fast Python Sampler for the von Mises-Fisher Distribution}, HAL-04004568v2, 2023.
\bibitem{kent1982} J. T. Kent, \textit{The Fisher-Bingham Distribution on the Sphere}, J. Royal Stat. Soc. B, 1982.
\bibitem{said2017} S. Said, L. Bombrun, Y. Berthoumieu, \textit{Riemannian Gaussian Distributions on the Space of Symmetric Positive Definite Matrices}, IEEE Trans. Inf. Theory, 2017.
\end{thebibliography}

\end{document}